% Options for packages loaded elsewhere
\PassOptionsToPackage{unicode}{hyperref}
\PassOptionsToPackage{hyphens}{url}
\PassOptionsToPackage{dvipsnames,svgnames,x11names}{xcolor}
\documentclass[
  letterpaper,
  10pt]{article}
\usepackage{xcolor}
\usepackage{graphicx}
\usepackage{float}
\usepackage[margin=0.8in]{geometry}
\usepackage{amsmath,amssymb}
\setcounter{secnumdepth}{-\maxdimen} % remove section numbering
\usepackage{iftex}
\ifPDFTeX
  \usepackage[T1]{fontenc}
  \usepackage[utf8]{inputenc}
  \usepackage{textcomp} % provide euro and other symbols
\else % if luatex or xetex
  \usepackage{unicode-math} % this also loads fontspec
  \defaultfontfeatures{Scale=MatchLowercase}
  \defaultfontfeatures[\rmfamily]{Ligatures=TeX,Scale=1}
\fi
\usepackage{lmodern}
\ifPDFTeX\else
  % xetex/luatex font selection
\fi
% Use upquote if available, for straight quotes in verbatim environments
\IfFileExists{upquote.sty}{\usepackage{upquote}}{}
\IfFileExists{microtype.sty}{% use microtype if available
  \usepackage[]{microtype}
  \UseMicrotypeSet[protrusion]{basicmath} % disable protrusion for tt fonts
}{}
\makeatletter
\@ifundefined{KOMAClassName}{% if non-KOMA class
  \IfFileExists{parskip.sty}{%
    \usepackage{parskip}
  }{% else
    \setlength{\parindent}{0pt}
    \setlength{\parskip}{6pt plus 2pt minus 1pt}}
}{% if KOMA class
  \KOMAoptions{parskip=half}}
\makeatother
\usepackage{color}
\usepackage{fancyvrb}
\newcommand{\VerbBar}{|}
\newcommand{\VERB}{\Verb[commandchars=\\\{\}]}
\DefineVerbatimEnvironment{Highlighting}{Verbatim}{commandchars=\\\{\}}
% Add ',fontsize=\small' for more characters per line
\newenvironment{Shaded}{}{}
\newcommand{\AlertTok}[1]{\textcolor[rgb]{1.00,0.00,0.00}{\textbf{#1}}}
\newcommand{\AnnotationTok}[1]{\textcolor[rgb]{0.38,0.63,0.69}{\textbf{\textit{#1}}}}
\newcommand{\AttributeTok}[1]{\textcolor[rgb]{0.49,0.56,0.16}{#1}}
\newcommand{\BaseNTok}[1]{\textcolor[rgb]{0.25,0.63,0.44}{#1}}
\newcommand{\BuiltInTok}[1]{\textcolor[rgb]{0.00,0.50,0.00}{#1}}
\newcommand{\CharTok}[1]{\textcolor[rgb]{0.25,0.44,0.63}{#1}}
\newcommand{\CommentTok}[1]{\textcolor[rgb]{0.38,0.63,0.69}{\textit{#1}}}
\newcommand{\CommentVarTok}[1]{\textcolor[rgb]{0.38,0.63,0.69}{\textbf{\textit{#1}}}}
\newcommand{\ConstantTok}[1]{\textcolor[rgb]{0.53,0.00,0.00}{#1}}
\newcommand{\ControlFlowTok}[1]{\textcolor[rgb]{0.00,0.44,0.13}{\textbf{#1}}}
\newcommand{\DataTypeTok}[1]{\textcolor[rgb]{0.56,0.13,0.00}{#1}}
\newcommand{\DecValTok}[1]{\textcolor[rgb]{0.25,0.63,0.44}{#1}}
\newcommand{\DocumentationTok}[1]{\textcolor[rgb]{0.73,0.13,0.13}{\textit{#1}}}
\newcommand{\ErrorTok}[1]{\textcolor[rgb]{1.00,0.00,0.00}{\textbf{#1}}}
\newcommand{\ExtensionTok}[1]{#1}
\newcommand{\FloatTok}[1]{\textcolor[rgb]{0.25,0.63,0.44}{#1}}
\newcommand{\FunctionTok}[1]{\textcolor[rgb]{0.02,0.16,0.49}{#1}}
\newcommand{\ImportTok}[1]{\textcolor[rgb]{0.00,0.50,0.00}{\textbf{#1}}}
\newcommand{\InformationTok}[1]{\textcolor[rgb]{0.38,0.63,0.69}{\textbf{\textit{#1}}}}
\newcommand{\KeywordTok}[1]{\textcolor[rgb]{0.00,0.44,0.13}{\textbf{#1}}}
\newcommand{\NormalTok}[1]{#1}
\newcommand{\OperatorTok}[1]{\textcolor[rgb]{0.40,0.40,0.40}{#1}}
\newcommand{\OtherTok}[1]{\textcolor[rgb]{0.00,0.44,0.13}{#1}}
\newcommand{\PreprocessorTok}[1]{\textcolor[rgb]{0.74,0.48,0.00}{#1}}
\newcommand{\RegionMarkerTok}[1]{#1}
\newcommand{\SpecialCharTok}[1]{\textcolor[rgb]{0.25,0.44,0.63}{#1}}
\newcommand{\SpecialStringTok}[1]{\textcolor[rgb]{0.73,0.40,0.53}{#1}}
\newcommand{\StringTok}[1]{\textcolor[rgb]{0.25,0.44,0.63}{#1}}
\newcommand{\VariableTok}[1]{\textcolor[rgb]{0.10,0.09,0.49}{#1}}
\newcommand{\VerbatimStringTok}[1]{\textcolor[rgb]{0.25,0.44,0.63}{#1}}
\newcommand{\WarningTok}[1]{\textcolor[rgb]{0.38,0.63,0.69}{\textbf{\textit{#1}}}}
\usepackage{longtable,booktabs,array}
\usepackage{calc} % for calculating minipage widths
% Correct order of tables after \paragraph or \subparagraph
\usepackage{etoolbox}
\makeatletter
\patchcmd\longtable{\par}{\if@noskipsec\mbox{}\fi\par}{}{}
\makeatother
% Allow footnotes in longtable head/foot
\IfFileExists{footnotehyper.sty}{\usepackage{footnotehyper}}{\usepackage{footnote}}
\makesavenoteenv{longtable}
\setlength{\emergencystretch}{3em} % prevent overfull lines
\providecommand{\tightlist}{%
  \setlength{\itemsep}{0pt}\setlength{\parskip}{0pt}}
\usepackage{bookmark}
\IfFileExists{xurl.sty}{\usepackage{xurl}}{} % add URL line breaks if available
\urlstyle{same}
\hypersetup{
  pdftitle={Separating Language Orchestration from Deterministic Control in Agentic Meal Recommendation},
  pdfauthor={Sai Teja Paladi, CSE 798 Independent Study},
  colorlinks=true,
  linkcolor={blue},
  filecolor={Maroon},
  citecolor={Blue},
  urlcolor={blue},
  pdfcreator={LaTeX via pandoc}}

\title{Separating Language Orchestration from Deterministic Control in
Agentic Meal Recommendation}
\author{Sai Teja Paladi, CSE 798 Independent Study}
\date{November 2026}

\begin{document}
\maketitle

\begin{abstract}

Large language models can translate natural-language requests into tool
calls, but delegating a complete recommendation process to a model
confounds language, tool-use, and recommendation errors. We study a
hybrid BEACON meal workflow in which models interpret requests while
deterministic components retain identity binding, Knapsack selection,
nutrition calculation, and validation. A hosted agent first reproduces
all six direct Knapsack plans, but is 28 times slower. On 30 requests
and five local models, an identity gate raises complete passes from
49/150 to 73/150 and reduces unsafe unclear-persona executions from ten
to zero without improving valid tool trajectories; a no-LLM pipeline
passes 30/30, showing that this request set does not justify model use.
In direct meal selection, Knapsack produces 6/6 safe plans with mean
Beacon quality 0.860; Granite produces 6/6 safe plans at 0.679, while
Llama produces only 3/6 safe plans despite 0.906 mean quality on that
safe subset. Finally, a 42-request study separates profile retrieval,
recommendation, nutrition analysis, and violation checking into four
tools. The deterministic baseline passes 38/42. Across 126 model-request
traces, prompt-only agents pass 62 cases and release all 18 injected
faulty plans; paired deterministic release enforcement raises passes to
80 and reduces unsafe releases to zero without changing model traces.
The results support using models only for language-dependent interaction
that outperforms a declared no-LLM baseline, while deterministic code
retains constraint-sensitive computation and release control.

\textbf{Keywords:} LLM agents, tool use, meal recommendation,
deterministic optimization, policy enforcement, clarification
\end{abstract}

\section{1. Introduction}\label{1-introduction}

Meal recommendation combines two different kinds of work.
Natural-language requests may be incomplete, ambiguous, or ask the
system to ignore stored restrictions, while profile retrieval,
constraint enforcement, and optimization should be reproducible. An LLM
can be useful for the first kind of work, but it is not automatically
the appropriate component for the second. Garimella, Srivastava, and
Sheth show this distinction clearly for exact arithmetic: deterministic
implementations achieve perfect accuracy with far lower latency than LLM
and agentic alternatives {[}3{]}. Their architectural decision framework
motivates using a model only where language understanding is useful and
retaining deterministic procedures for exact work.

This paper applies that principle to the BEACON meal-recommendation
setting. BEACON studies the tradeoff between nutrition and convenience
for configurable, long-horizon meal recommendations {[}11, 12{]}. We do
not replace or modify its Knapsack recommender. Instead, we investigate
whether an LLM can reliably orchestrate the workflow around it and what
happens when the user\textquotesingle s persona is missing, unknown, or
ambiguous.

The study asks five research questions:

\begin{itemize}
\tightlist
\item
  \textbf{RQ1:} Can one hosted LLM complete the required tool trajectory
  while preserving the exact output of direct deterministic Knapsack
  execution?
\item
  \textbf{RQ2:} How robust is prompt-only orchestration across
  natural-language, restriction-override, missing, unknown, and
  ambiguous-persona requests?
\item
  \textbf{RQ3:} Can a deterministic identity gate improve clarification
  safety without changing the recommender?
\item
  \textbf{RQ4:} How do direct LLM-selected meal plans compare with
  deterministic Knapsack plans on constraint adherence, quality,
  nutrition fit, variety, and latency?
\item
  \textbf{RQ5:} Can responsibility-based tools support multiple user
  intents, and does deterministic validation prevent a corrupted plan
  from reaching the response boundary?
\end{itemize}

The paper makes five bounded contributions:

\begin{enumerate}
\def\labelenumi{\arabic{enumi}.}
\tightlist
\item
  a controlled separation between language-model orchestration and
  deterministic meal selection in the BEACON system.
\item
  an identity-gated hybrid architecture that resolves a persona outside
  the model, constrains identity-dependent tool calls, and renders
  verified results without a final generative step.
\item
  a trace-based, five-model evaluation that measures exact plan
  fidelity, tool order, clarification, unsafe execution, failures, and
  latency without an LLM judge.
\item
  a same-persona, same-catalog comparison between deterministic Knapsack
  and two local LLM-only selectors, scored after generation by the same
  Beacon evaluators.
\item
  a four-tool, multi-intent extension with deterministic nutrition
  analysis, violation checking, injected-fault evaluation, and a matched
  release-gate ablation across three local models.
\end{enumerate}

This is not a claim of a new optimization algorithm, a new language
model, or a clinically validated recommender. It is a controlled systems
study of agentic orchestration around an existing deterministic
recommender.

\section{2. Related Work}\label{2-related-work}

\subsection{2.1 Reasoning, acting, and recommendation
agents}\label{21-reasoning-acting-and-recommendation-agents}

ReAct formalizes an agent pattern in which a language model interleaves
reasoning, actions, and observations from an external environment
{[}1{]}. Our implementation follows the observable part of this pattern:
the model receives tool schemas, chooses tool names and arguments,
observes structured results, and decides whether to call another tool or
stop. The presence of a particular agent framework is not what makes
this loop agentic.

RecMind applies an autonomous LLM agent to recommendation using
planning, personalized memory, world knowledge, and tools {[}2{]}. It
demonstrates that LLMs can add flexible language understanding and
planning to recommendation tasks. Our study asks a different and
narrower question. The LLM does not rank meals; it must correctly invoke
an existing recommender. This boundary allows an orchestration failure
to be separated from a meal-selection failure.

\subsection{2.2 Choosing between deterministic and agentic
computation}\label{22-choosing-between-deterministic-and-agentic-computation}

\emph{Doing Less with More} compares native functions, service
architectures, direct LLM responses, and tool-using agents on 1,000
exact expressions {[}3{]}. Its central lesson is not that agents are
universally ineffective, but that architecture should follow task
properties. Exact known procedures should remain deterministic, while
models can be reserved for linguistic or uncertain steps. Our
architecture transfers this idea from arithmetic to meal-agent workflow
control: Knapsack selection and identity resolution remain
deterministic; request interpretation and valid tool sequencing remain
model-driven.

\subsection{2.3 Evaluating tool-using
agents}\label{23-evaluating-tool-using-agents}

Tool-use evaluation requires more than judging final prose.
\(\tau\)-bench evaluates tool-agent-user interaction using domain
policies and final environment state {[}4{]}. ToolSandbox emphasizes
state dependencies, canonicalization, and insufficient-information cases
{[}5{]}. The Berkeley Function-Calling Leaderboard also evaluates tool
selection, arguments, multi-turn execution, and irrelevant or
hallucinated calls {[}6{]}. Following this line of work, our evaluator
checks the executed persona, tool order, Knapsack action, and exact
returned plan rather than using an LLM to rate the response text.

The Agentic AI Test Checklist collects 110 implementation-oriented
checks for terminal state, dependency ordering, tool allowlists, runtime
enforcement, and related risks {[}15{]}. It is not a meal-recommendation
benchmark, so we do not claim a score over the full checklist. We adapt
selected checks into task-specific oracles: all postconditions must hold
together, dependent tools must follow successful prerequisites, the tool
set and turn count are bounded, and plan release is enforced by the
application rather than only by a prompt.

\subsection{2.4 Clarification and runtime policy
enforcement}\label{24-clarification-and-runtime-policy-enforcement}

Clarification is a separate capability from tool invocation. Testoni and
Fernández show that model uncertainty does not necessarily match human
clarification behavior {[}7{]}. Zhang and Choi likewise frame deciding
when to ask as an explicit stage in ambiguity resolution {[}8{]}. These
findings motivate treating missing or ambiguous identity as an
observable test condition rather than assuming that a system prompt will
reliably produce a question.

AgentDojo shows that tool-using agents may fail tasks and security
properties even without or under prompt-injection attacks {[}9{]}. CaMeL
proposes placing a protective system layer around the model and
enforcing capabilities through explicit control and data flows {[}10{]}.
Our policy gate is substantially smaller and does not inherit
CaMeL\textquotesingle s formal guarantees. It adopts the relevant design
principle: a consequential identity-binding decision should be enforced
by the application rather than delegated solely to prompt compliance.
The current prototype performs identity resolution, not authentication:
it does not verify that the requester is entitled to access the named
profile.

\subsection{2.5 BEACON and meal
recommendation}\label{25-beacon-and-meal-recommendation}

BEACON formulates meal planning as a combination problem involving food
roles, user constraints, convenience, nutrition, and longer-term
considerations {[}11, 12{]}. Its prototype includes deterministic
recommendation methods and quality measures. Our study reuses its
personas, local food sources, Knapsack implementation, and evaluators.
We hold these components fixed so the experiment measures orchestration
rather than introducing a competing meal algorithm.

\subsection{2.6 MCP and LangChain}\label{26-mcp-and-langchain}

The Model Context Protocol standardizes how AI applications discover and
invoke external tools and resources {[}13{]}. LangChain provides an
agent harness that can implement a model-tool loop {[}14{]}. Neither
changes the scientific distinction between model-controlled
orchestration and deterministic selection. The present benchmark
therefore uses a small native OpenAI-compatible function-calling loop,
which makes every request, response, and tool trace directly auditable.
An MCP or LangChain implementation would be a useful portability or
framework-overhead comparison, but it is not the novelty claimed here.

\section{3. System Design}\label{3-system-design}

\subsection{3.1 Fixed recommendation
component}\label{31-fixed-recommendation-component}

For a stored persona \(p\), the BEACON Knapsack implementation produces
a one-day plan:

\[
K(p, C, d) \rightarrow M,
\]

where \(C\) is the fixed local catalog, \(d\) is the fixed planning
date, and \(M\) contains breakfast, lunch, and dinner role-item
assignments. The same persona, catalog, date, and backend version are
used in all conditions.

The language model never supplies meal IDs to Knapsack and never
replaces the selected plan. Consequently, if two conditions successfully
execute Knapsack for the same persona, their expected role-item
assignments are identical.

\subsection{3.2 Agent tools}\label{32-agent-tools}

The agent receives three logical tools:

\begin{table}[H]
\centering
\scriptsize
\resizebox{\textwidth}{!}{%
\begin{tabular}{@{}lll@{}}
\toprule\noalign{}
Tool & Purpose & State effect \\
\midrule\noalign{}
\texttt{list\_personas()} & Return stored persona identifiers and names.
& Read only \\
\texttt{get\_persona(persona\_id)} & Retrieve one stored profile. &
Records retrieved identity \\
\texttt{run\_knapsack\_for\_persona(persona\_id)} & Execute BEACON
Knapsack. & Produces the meal plan \\
\bottomrule\noalign{}
\end{tabular}}
\end{table}

For valid requests, the required trajectory is:

\begin{Shaded}
\begin{Highlighting}[]
\NormalTok{user request}
\NormalTok{    → model calls list\_personas}
\NormalTok{    → model calls get\_persona for the intended persona}
\NormalTok{    → model calls run\_knapsack\_for\_persona for the same persona}
\NormalTok{    → verified deterministic plan}
\end{Highlighting}
\end{Shaded}

The model genuinely chooses each tool and its arguments. A wrong tool,
wrong persona, wrong order, missing call, or repeated execution is
preserved in the trace and scored as a failure.

\subsection{3.3 Prompt-only
architecture}\label{33-prompt-only-architecture}

In the prompt-only condition, the system prompt tells the model to
identify the persona, retrieve the profile, preserve stored
restrictions, ask for clarification when necessary, and invoke Knapsack.
These instructions are not enforced outside the model. The model can
therefore choose an unrelated persona, skip a step, provide the full
name where an internal identifier is required, or run Knapsack despite
insufficient information.

\subsection{3.4 Identity-gated
architecture}\label{34-identity-gated-architecture}

The proposed architecture adds four deterministic controls:

\begin{enumerate}
\def\labelenumi{\arabic{enumi}.}
\tightlist
\item
  \textbf{Identity resolution.} Before model execution, application code
  performs a case- and punctuation-insensitive match against one
  explicit full name or unique first name in the stored persona set.
\item
  \textbf{Safe clarification.} Missing, unknown, or multiple persona
  matches return a deterministic clarification. The model and
  identity-dependent tools are not called.
\item
  \textbf{Argument restriction.} For a resolved request,
  \texttt{get\_persona} and \texttt{run\_knapsack\_for\_persona} are
  locked to the resolved identifier. Knapsack also requires the list and
  matching-profile steps to have succeeded.
\item
  \textbf{Trusted rendering.} After successful Knapsack execution, the
  application returns a short confirmation from trusted state. It does
  not ask the model to restate nutrition facts or invent an explanation.
\end{enumerate}

This gate does not choose meals and does not complete the valid tool
sequence for the model. It converts incorrect model actions into
explicit failures while preventing them from becoming wrong-person
recommendations.

\subsection{3.5 Responsibility-based tool
extension}\label{35-responsibility-based-tool-extension}

The initial three tools expose implementation steps rather than all of
the responsibilities discussed for a meal agent. The extension uses four
tools with distinct inputs, outputs, and failure boundaries:

\begin{table}[H]
\centering
\scriptsize
\resizebox{\textwidth}{!}{%
\begin{tabular}{@{}lll@{}}
\toprule\noalign{}
Tool & Responsibility & Enforced prerequisite \\
\midrule\noalign{}
\texttt{fetch\_persona\_profile} & Resolve one ID, full name, or unique
first name and return the stored fictional profile. &
Application-resolved persona must match. \\
\texttt{recommend\_meal\_plan} & Run the unchanged deterministic Beacon
Knapsack method. & Matching profile was fetched. \\
\texttt{calculate\_nutrition\_scores} & Return nutrient totals, goal
fit, and FSA/WHO-style proxies when requested. & A current plan
exists. \\
\texttt{check\_plan\_violations} & Detect structural and stored
dietary-condition violations. & A current plan exists. \\
\bottomrule\noalign{}
\end{tabular}}
\end{table}

Profile selection and retrieval are combined because they constitute one
read-only database access responsibility. Recommendation, nutrition
analysis, and validation remain separate because they answer different
intents and fail independently. Plans are held in trusted session state,
so the model does not copy a large plan into the nutrition or validation
tools.

For a resolved request, the model still chooses which tools to invoke.
The application binds all calls to the resolved profile and enforces
dependencies. In the policy-gated condition, it also refuses to release
a plan unless the violation checker has passed. A prompt-only ablation
disables only this final release rule; it does not change the prompt,
schemas, model, or recommender.

The FSA-style value is the mean of available per-item traffic-light
proxies using sugar, salt converted from sodium, fat, and saturated fat
(range 4-\/-12; lower is better). The WHO-style value is the mean of
available custom per-item range checks over seven nutrients (range
0-\/-7; higher is better). Coverage is reported and missing values are
not imputed. Neither proxy is an official label, clinical score, or
individualized advice.

For example, Earl\textquotesingle s deterministic plan contains nine
selected catalog entries, but only five have the complete fields
required by both proxies. The tool therefore returns FSA-style 9.8 and
WHO-style 3.8 with coverage 5/9, rather than silently filling the four
missing entries. Its nutrition-goal fit is 0.506: the calorie component
is 1.000, while carbohydrate, protein, and fiber fits are 0.307, 0.530,
and 0.187. Nutrient totals sum available catalog values and must be read
together with the missing-data fields.

\section{4. Experimental Method}\label{4-experimental-method}

\subsection{4.1 Experiment 1: direct versus agentic integration
check}\label{41-experiment-1-direct-versus-agentic-integration-check}

Experiment 1 is an integration and tool-compliance check for RQ1 using
the six fictional BEACON personas. The direct condition receives a
selected persona identifier and invokes Knapsack. The agentic condition
receives a natural-language request, uses hosted
\texttt{openai/gpt-oss-20b} through Groq at temperature zero, and must
complete the three-tool trajectory. Both conditions use the same
backend, catalog, fixed date, and reference plans. Because Knapsack is
deterministic, identical meal output is expected whenever the model
supplies the correct persona identifier; the nontrivial measurement is
whether this one hosted model completes the required tool trajectory.

The metrics are requested-role coverage, exact role-item fidelity to
direct Knapsack, required tool-sequence compliance,
BEACON\textquotesingle s original goodness, BEACON\textquotesingle s
newer quality score, post-hoc nutrition-goal fit, failures, retries, and
end-to-end latency. Nutrition fit is evaluated after generation; the
tested Knapsack implementation does not optimize the numeric calorie,
carbohydrate, protein, and fiber targets directly.

\subsection{4.2 Experiment 2: request
robustness}\label{42-experiment-2-request-robustness}

Experiment 2 tests RQ2 and RQ3 with 30 fixed requests over the same six
stored personas. Twenty-four requests require recommendation: six each
use an exact full name, a paraphrase, a unique first name, or a
restriction-override request that asks the system to ignore saved
restrictions. Six requests require clarification because identity is
unknown, missing, or ambiguous.

Five local Ollama models are evaluated:

\begin{itemize}
\tightlist
\item
  \texttt{llama3.2:1b};
\item
  \texttt{llama3.2:3b};
\item
  \texttt{phi4-mini:3.8b};
\item
  \texttt{qwen3.5:4b};
\item
  \texttt{granite4:3b}.
\end{itemize}

Each model receives every request in both the prompt-only and
identity-gated architectures, producing 300 case-level executions. A
no-LLM pipeline also applies the same name resolver and then runs
Knapsack directly. Runs are sequential to avoid local memory contention.
Temperature is zero, the turn limit is eight, and the same tool schemas,
backend, reference plans, and scoring rules are used.

\subsection{4.3 Experiment 3: deterministic versus LLM-only
selection}\label{43-experiment-3-deterministic-versus-llm-only-selection}

Experiment 3 tests RQ4 using the same six personas and local Beacon
catalog. The deterministic condition runs Knapsack. In each LLM-only
condition, the model receives the stored profile and all 124 unique
candidates eligible for the requested main-course, side, or beverage
roles (145 role memberships) in a compact representation containing
candidate ID, name, role, dietary flags, and available calories,
carbohydrates, protein, and fiber. The LLM selects the nine catalog
items directly. It cannot call Knapsack or another recommendation tool,
and candidates are not pre-filtered for the persona\textquotesingle s
dietary conditions.

To measure meal choice rather than malformed ID copying, Ollama
structured output constrains each field to an ID belonging to the
requested role and catalog. This schema performs no ranking and applies
no persona-specific dietary filter. \texttt{granite4:3b} and
\texttt{llama3.2:3b} each make one call per persona at temperature zero.
After selection, deterministic validation checks stored dietary
conditions, and the same Beacon evaluators measure goodness, overall
quality, nutrition fit, repeated foods, exact Knapsack match, tokens,
and latency.

\subsection{4.4 Experiment 4: meaningful tools and fault
containment}\label{44-experiment-4-meaningful-tools-and-fault-containment}

Experiment 4 evaluates the four responsibility-based tools on 42 fixed
non-sensitive requests, six in each of seven categories: profile-only,
recommendation, nutrition analysis, restriction override, indirect
intent, unclear identity, and injected fault. Unclear identities are
handled before model execution by the same deterministic resolver. For
every other case, the model must select an intent-appropriate sequence
rather than replay one fixed trajectory.

The injected fault removes the first required role from a completed plan
after Knapsack and before validation. This isolates whether the checker
detects the fault and whether the release boundary enforces that result.
It does not test whether the model can notice the corruption from text
because the model never receives or edits the plan directly.

Three local Ollama models are used: \texttt{granite4:3b},
\texttt{qwen3.5:4b}, and \texttt{llama3.2:3b}. Each model receives all
42 requests once with the release gate disabled. To isolate the gate
from model nondeterminism, each of the 126 model-request traces is then
replayed through the deterministic gate without another model call. A
deterministic keyword pipeline runs the same 42 cases as a no-LLM
baseline. All runs use the same dataset, four tool schemas, backend,
Knapsack implementation, temperature zero, and eight-turn limit.

\subsection{4.5 Scoring}\label{45-scoring}

A valid request passes only if it:

\begin{enumerate}
\def\labelenumi{\arabic{enumi}.}
\tightlist
\item
  retrieves the intended profile;
\item
  invokes Knapsack exactly once for that persona;
\item
  follows \texttt{list\ →\ profile\ →\ Knapsack};
\item
  exactly matches the direct Knapsack role-item assignments.
\end{enumerate}

An unclear request passes only if it returns a usable clarification and
does not run Knapsack. We separately count unsafe Knapsack executions,
policy-blocked calls, API or loop failures, failed tools, and wall-clock
latency. No generated text is graded for recommendation quality by
another model.

An LLM-only plan is safe and valid only if it fills every requested meal
role with a catalog member and violates none of the stored diabetes,
vegetarian, vegan, or gluten-free conditions. Quality metrics are
computed after generation and do not influence the
model\textquotesingle s selection.

For Experiment 4, complete-case success requires all applicable checks
to pass: persona routing, exact tool membership with dependency-valid
order, expected terminal action, correct nutrition-tool selection, exact
Knapsack fidelity, zero failed tools, and zero API or loop errors. A
fault case also requires detection and blocked release. We separately
report unsafe releases and success by request category. The metric
dictionary gives each denominator and a worked example.

\subsection{4.6 Reproducibility}\label{46-reproducibility}

The benchmark records the dataset hash, prompt hash, source commit,
hashes of the benchmark and policy-gate implementations, model,
temperature, complete tool trace, application errors, and latency. All
profiles are fictional and all food sources are local to the compatible
BEACON checkout.

\section{5. Results}\label{5-results}

\subsection{5.1 Direct versus agentic integration
check}\label{51-direct-versus-agentic-integration-check}

\begin{table}[H]
\centering
\scriptsize
\resizebox{\textwidth}{!}{%
\begin{tabular}{@{}lrrrrrrrr@{}}
\toprule\noalign{}
Condition & Completed & Role coverage & Exact fidelity & Required
sequence & Goodness & BEACON quality & Nutrition fit & Mean latency \\
\midrule\noalign{}
Direct Knapsack & 6/6 & 100\% & 100\% & N/A & 1.000 & 0.860 & 0.549 &
202.4 ms \\
Agentic Knapsack & 6/6 & 100\% & 100\% & 100\% & 1.000 & 0.860 & 0.549 &
5706.4 ms \\
\bottomrule\noalign{}
\end{tabular}}
\end{table}

The hosted agent completes the required tool sequence and reproduces the
deterministic output for all six personas. This is evidence of
integration and tool-use compliance for one model, not evidence that the
LLM discovers a better recommendation. Meal-quality metrics are
identical because the same Knapsack method selects the food. The agentic
path is approximately 28 times slower when provider retry waits are
included. Excluding 23.8 seconds of recorded rate-limit waits, its mean
latency is 1747.2 ms, still about 8.6 times the direct baseline.

\subsection{5.2 Prompt-only, identity-gated, and no-LLM
robustness}\label{52-prompt-only-identity-gated-and-no-llm-robustness}

\begin{table}[H]
\centering
\scriptsize
\resizebox{\textwidth}{!}{%
\begin{tabular}{@{}llrrrrr@{}}
\toprule\noalign{}
Model & Architecture & Overall & Valid requests & Unclear requests &
Unsafe executions & Mean latency \\
\midrule\noalign{}
None & No-LLM pipeline & \textbf{30/30} & \textbf{24/24} & \textbf{6/6}
& \textbf{0} & \textbf{146.3 ms} \\
Granite 4 3B & Prompt-only & 22/30 & 22/24 & 0/6 & 6 & 9815.1 ms \\
Granite 4 3B & Identity-gated & \textbf{28/30} & 22/24 & \textbf{6/6} &
\textbf{0} & \textbf{3154.7 ms} \\
Llama 3.2 1B & Prompt-only & 0/30 & 0/24 & 0/6 & 0 & 3596.2 ms \\
Llama 3.2 1B & Identity-gated & \textbf{6/30} & 0/24 & \textbf{6/6} & 0
& \textbf{2150.2 ms} \\
Llama 3.2 3B & Prompt-only & 5/30 & 5/24 & 0/6 & 3 & 3868.7 ms \\
Llama 3.2 3B & Identity-gated & \textbf{11/30} & 5/24 & \textbf{6/6} &
\textbf{0} & \textbf{2809.4 ms} \\
Phi-4 Mini 3.8B & Prompt-only & 6/30 & 0/24 & 6/6 & 0 & \textbf{1992.8
ms} \\
Phi-4 Mini 3.8B & Identity-gated & 6/30 & 0/24 & 6/6 & 0 & 2001.4 ms \\
Qwen 3.5 4B & Prompt-only & 16/30 & 16/24 & 0/6 & 1 & 15351.0 ms \\
Qwen 3.5 4B & Identity-gated & \textbf{22/30} & 16/24 & \textbf{6/6} &
\textbf{0} & \textbf{8953.8 ms} \\
\bottomrule\noalign{}
\end{tabular}}
\end{table}

Aggregating the matched runs gives the following result:

\begin{table}[H]
\centering
\scriptsize
\resizebox{\textwidth}{!}{%
\begin{tabular}{@{}lrr@{}}
\toprule\noalign{}
Measure & Prompt-only & Identity-gated \\
\midrule\noalign{}
Complete-case passes & 49/150 & \textbf{73/150} \\
Valid workflow passes & 43/120 & 43/120 \\
Unclear-request passes & 6/30 & \textbf{30/30} \\
Unsafe unclear-case Knapsack executions & 10 & \textbf{0} \\
Mean latency per case & 6924.8 ms & \textbf{3813.9 ms} \\
\bottomrule\noalign{}
\end{tabular}}
\end{table}

The policy gate improves all 24 unclear cases that were not already
complete passes and introduces no matched pass-to-fail regressions. The
44.9\% mean latency reduction comes from bypassing model inference for
unclear identity and omitting an unnecessary final model turn after
verified Knapsack execution.

The no-LLM pipeline is both faster and more accurate on this dataset.
Every valid request contains an exact full name or unique first name,
every unclear case is covered by the resolver\textquotesingle s rules,
and the remaining workflow is fixed. Its 30/30 result is therefore close
to true by construction. It establishes an important negative result:
these 30 cases do not show that an LLM is needed. Requests involving
typos, nicknames, contextual references, or free-form preferences are
needed to test added language-understanding value.

The 43/120 valid-workflow total is identical because the same valid
requests fail in matched prompt-only and gated runs. Across the five
prompt-only runs, 43 cases pass, 41 fail in backend tool calls because
the model supplies an invalid persona argument, and 36 never make a
successful Knapsack call. In the gated runs, the 41 backend errors
become 41 explicit policy blocks; the other 36 still lack a successful
call. The gate contains invalid arguments but does not repair the
model\textquotesingle s action or force the sequence.

\subsection{5.3 Deterministic versus LLM-only
selection}\label{53-deterministic-versus-llm-only-selection}

\begin{table}[H]
\centering
\scriptsize
\resizebox{\textwidth}{!}{%
\begin{tabular}{@{}lrrrrrrrrr@{}}
\toprule\noalign{}
Selector & Safe plans & Hard violations & Exact Knapsack matches &
Goodness & Quality, all cases & Quality, safe only & Nutrition fit &
Repeated foods & Latency \\
\midrule\noalign{}
Beacon Knapsack & \textbf{6/6} & \textbf{0} & 6/6 reference &
\textbf{1.000} & \textbf{0.860} & 0.860 & 0.549 & \textbf{0.0} &
\textbf{183.5 ms} \\
Granite 4 3B LLM-only & \textbf{6/6} & \textbf{0} & 0/6 & \textbf{1.000}
& 0.679 & 0.679 & 0.479 & 2.0 & 4338.9 ms \\
Llama 3.2 3B LLM-only & 3/6 & 4 & 0/6 & 0.981 & 0.453 & \textbf{0.906} &
\textbf{0.563} & \textbf{0.0} & 4035.3 ms \\
\bottomrule\noalign{}
\end{tabular}}
\end{table}

Granite returns six dietary-adherent plans, but its mean Beacon quality
is 0.181 below Knapsack, nutrition fit is lower, and it repeats two
foods per plan on average. Llama avoids repeated foods and slightly
exceeds Knapsack\textquotesingle s mean nutrition fit, but it violates
stored dietary conditions for Jessica, Becca, and Dylan. The four
violations are two non-low-sugar choices for a persona with diabetes and
one meat-containing choice for each of two vegetarians. Beacon assigns
zero overall quality to those unsafe plans, reducing
Llama\textquotesingle s mean quality to 0.453. Its three safe plans
average 0.906, above Knapsack\textquotesingle s 0.860. This conditional
result is based on only three selected cases and must be read together
with Llama\textquotesingle s 3/6 safety rate; it prevents a claim that
Knapsack has the highest conditional meal quality.

Neither model exactly reproduces a Knapsack plan. This demonstrates that
the models are genuinely selecting different meals; exact disagreement
alone is not an error. The safety and quality metrics show why the
difference matters. Both LLMs are more than 20 times slower than direct
Knapsack and consume about 12,900 prompt-plus-completion tokens across
six personas, whereas Knapsack uses no model tokens.

\subsection{5.4 Meaningful-tool and fault-containment
result}\label{54-meaningful-tool-and-fault-containment-result}

\begin{table}[H]
\centering
\scriptsize
\resizebox{\textwidth}{!}{%
\begin{tabular}{@{}lrrrrr@{}}
\toprule\noalign{}
Model or baseline & Prompt-only & Policy-gated paired replay & Valid
tool path & Unsafe releases, prompt → gated & Mean latency \\
\midrule\noalign{}
Deterministic rule pipeline & 38/42 & N/A & 90.5\% & 0 → 0 & 264.9 ms \\
Granite 4 3B & 24/42 & \textbf{30/42} & 71.4\% & 6 → \textbf{0} & 4573.8
ms \\
Llama 3.2 3B & 19/42 & \textbf{25/42} & 61.9\% & 6 → \textbf{0} & 3630.6
ms \\
Qwen 3.5 4B & 19/42 & \textbf{25/42} & 59.5\% & 6 → \textbf{0} & 7843.7
ms \\
\bottomrule\noalign{}
\end{tabular}}
\end{table}

Across the same 126 model-request traces, the release gate increases
complete passes from 62 to 80 and reduces unsafe release of injected
faulty plans from 18 to zero. Every increase is a fault case whose model
called the violation checker and detected the missing role but whose
prompt-only application would still release the plan. Because each trace
is replayed rather than regenerated, the gate cannot change identity
routing, tool selection, Knapsack output, model failures, or latency.

The broader workflow remains difficult. Granite has 100\% Knapsack
fidelity over the 24 ordinary recommendation-required cases, but it
completes only one of six plain recommendation cases and none of six
nutrition-analysis cases because it usually omits the final violation
check. Llama and Qwen have only 41.7\% and 29.2\% Knapsack fidelity
respectively because they often stop before producing a plan or choose
an invalid workflow. Llama records 27 failed local calls and two
turn-limit failures; the other two models have no API or loop failures.

Category results also show why one overall score is insufficient. All
three agents handle six of six profile-only requests, whereas the
keyword baseline handles three. Granite passes all restriction-override
requests, while Llama and Qwen pass none. Nutrition-analysis success is
0/6 for Granite, 2/6 for Llama, and 1/6 for Qwen. All unclear-identity
cases pass because the application resolver clarifies before model
execution; this is a system result, not model evidence. The
deterministic pipeline\textquotesingle s 38/42 remains the strongest
overall result and is 13.7-\/-29.6 times faster than the three agent
runs.

\subsection{5.5 Model capability versus architecture
safety}\label{55-model-capability-versus-architecture-safety}

The valid-workflow total is exactly 43/120 in both architectures. This
is the most important interpretation boundary in the study: the gate
does not improve the model\textquotesingle s ability to select tools.
Granite and Qwen remain substantially more capable than the smaller
models on valid trajectories. The gate instead changes what the system
permits when the model is wrong or the request is incomplete.

Phi-4 Mini illustrates the distinction. It safely clarifies all six
unclear requests even without the gate but completes none of the 24
valid workflows. Granite illustrates the opposite failure: it completes
22/24 valid workflows but runs Knapsack on all six unclear prompt-only
requests. General tool-use capability and safe clarification are
therefore empirically separable.

\section{6. Discussion}\label{6-discussion}

\subsection{6.1 Answering the research
questions}\label{61-answering-the-research-questions}

\textbf{RQ1:} One hosted model completes all six required tool
trajectories and preserves the fixed Knapsack output. This is a small
integration result. It adds substantial latency and does not test
whether language understanding is useful.

\textbf{RQ2:} Prompt-only robustness varies widely by model. No tested
model combines high valid-workflow completion with consistently safe
unclear-persona handling. The strongest prompt-only model completes
22/24 valid cases but executes the recommender in all six cases where
clarification is required.

\textbf{RQ3:} The deterministic gate produces complete safe
clarification for all 30 unclear model-request pairs, eliminates ten
unsafe executions, and reduces mean latency without changing
valid-workflow performance. It improves the system, not the underlying
models. The no-LLM baseline\textquotesingle s 30/30 result further shows
that the present request set does not require a model.

\textbf{RQ4:} In this preliminary controlled comparison, Granite
preserves dietary constraints but has lower all-case quality, nutrition
fit, and variety than Knapsack. Llama violates hard conditions in half
of the personas, although its three safe plans have higher conditional
quality. Knapsack is faster, token-free, and safe on all six cases. The
small sample and conditional result do not support a general ranking of
selector quality.

\textbf{RQ5:} Responsibility-based tools make profile retrieval,
recommendation, nutrition analysis, and validation separately
observable, but none of the three models matches the 38/42 deterministic
baseline. The paired release gate turns all 18 detected injected faults
into blocked outcomes without changing a tool trace. It improves system
safety, not language understanding or tool-use capability.

\subsection{6.2 What is novel}\label{62-what-is-novel}

The novelty is the combination of a task-specific architecture and
controlled evaluation:

\begin{itemize}
\tightlist
\item
  it isolates natural-language orchestration from deterministic meal
  selection;
\item
  it converts persona resolution into an explicit application-level
  control decision instead of relying only on a system prompt;
\item
  it evaluates the same agent with and without that gate across five
  local models and identical requests;
\item
  it scores executed state---identity, tool order, Knapsack action, and
  exact plan fidelity---while separately reporting safety and model
  capability;
\item
  it compares deterministic and direct LLM selection on the same
  profiles, catalog, roles, and post-hoc evaluators instead of comparing
  systems on incompatible inputs.
\item
  it evaluates four responsibility-based tools on multiple intents and
  isolates the release gate\textquotesingle s effect with trace-paired
  fault replay rather than separate model samples.
\end{itemize}

The work does not claim that deterministic identity resolution is itself
a new idea. Its contribution is applying and measuring that principle in
a reproducible BEACON meal-agent workflow, exposing a concrete tradeoff
that a final-answer accuracy metric would hide: a model can be good at
valid tool trajectories and simultaneously unsafe when identity is
unclear.

\subsection{6.3 Architecture lesson from direct LLM
selection}\label{63-architecture-lesson-from-direct-llm-selection}

The direct-selection comparison clarifies the architecture tradeoff. The
LLM-only selectors are not merely slower alternatives that produce the
same plans. They make different tradeoffs: Granite repeatedly chooses
foods despite being safe, while Llama obtains a favorable numeric
nutrition-fit mean partly through plans that violate stored dietary
conditions. A single aggregate metric would hide both failures.

The result supports a hybrid design in which deterministic code first
enforces hard constraints and produces feasible candidates. An LLM may
still add value by reranking that bounded set using soft preferences or
by explaining tradeoffs. This is different from allowing the model to
select from the unrestricted catalog and hoping a prompt will preserve
every condition.

\section{7. Limitations and Ethical
Scope}\label{7-limitations-and-ethical-scope}

The evaluation uses six fictional personas and 30 handcrafted requests.
Exact first-name matching is deliberately conservative and does not
resolve typos, nicknames, pronouns, or conversational identity. The
study uses one local machine and one run per model-request pair at
temperature zero; it does not estimate variance across sampling,
hardware, or model versions. The models are small local models, so
results should not be generalized to frontier hosted systems.

The responsibility-based extension adds 42 handcrafted cases but uses
the same six personas, one prompt per category and persona, and three
small local models. Its unclear-identity outcomes come from the
pre-model resolver and do not test model clarification. Its fault set
contains one synthetic corruption---removing the first required meal
role---so zero unsafe releases demonstrates enforcement for that fault,
not complete safety. The paired replay isolates the
gate\textquotesingle s effect, but it does not estimate variation across
repeated model generations.

The gate\textquotesingle s six unclear cases were handcrafted for the
resolver\textquotesingle s fixed rules, so its perfect clarification
result is an invariant on this test set rather than evidence of
generalization. Name matching is identity resolution, not
authentication: anyone can type another persona\textquotesingle s name.
A stronger prompt-only baseline with few-shot examples or an explicit
\texttt{ask\_clarification} tool was not tested. Those baselines,
repeated trials, and language-dependent requests are required before
claiming that the gate is preferable to the strongest prompting or that
an LLM adds value over the no-LLM pipeline.

The LLM-only experiment uses role- and catalog-constrained JSON
generation. It therefore measures meal choice after eliminating
malformed role/ID output; it does not measure fully unconstrained
generation. It does not filter candidates by persona, so the observed
dietary violations remain model decisions. Only two models are included
in this recommendation-source comparison.

Meal plans are evaluated using existing BEACON metrics and deterministic
fidelity. The study does not establish that plans are clinically
appropriate, that users will follow them, or that the policy gate
provides complete security. Nutrition-goal fit is post-hoc and must not
be described as an optimized constraint. No real personal or protected
health information is used. Because Knapsack and the BEACON evaluators
come from the same system, the quality measures are not an independent
human judgment and may align with the deterministic
method\textquotesingle s design objectives.

The FSA-style and WHO-style outputs are application-defined per-serving
proxies, not official labels. In the worked plan, only five of nine
selected entries have complete inputs, and nutrient totals can be
partial when catalog fields are missing. Coverage and missing-data
fields must accompany every reported value.

The aggregate 49/150 and 73/150 observations contain repeated prompts
across models and should not be treated as 150 independent human
samples. The current result is a controlled systems comparison, not a
user study or clinical trial.

\section{8. Next Experiment: Constrained Hybrid
Reranking}\label{8-next-experiment-constrained-hybrid-reranking}

The next study should add a hybrid selection condition while retaining
the two completed baselines. It should hold personas, catalog, candidate
visibility, and evaluation metrics fixed and compare:

\begin{enumerate}
\def\labelenumi{\arabic{enumi}.}
\tightlist
\item
  \textbf{Deterministic Knapsack:} the existing reference selector;
\item
  \textbf{LLM-only selector:} the completed direct-selection condition;
\item
  \textbf{Constrained hybrid reranker:} deterministic code produces
  several feasible candidate plans, the LLM reranks them for stated soft
  preferences, and a deterministic validator rejects any invalid output.
\end{enumerate}

Run each stochastic condition repeatedly and expand the persona/request
set. Primary measures should include hard-constraint violations, BEACON
goodness, quality and nutrition-goal fit, preference satisfaction, exact
consistency over repeated runs, explanation faithfulness, latency, and
token cost. This three-way experiment would answer whether
language-model judgment adds recommendation value. It must remain
separate from the orchestration experiment so a ranking change is not
confused with a routing failure.

\section{9. Conclusion}\label{9-conclusion}

This study demonstrates that an LLM can correctly orchestrate an
existing deterministic meal recommender, but success on clear requests
does not ensure safe behavior when identity is insufficient. Prompt-only
agents either fail valid workflows, fail clarification, or both. A small
deterministic policy gate eliminates observed wrong-person Knapsack
executions and reduces latency while preserving the
models\textquotesingle{} valid-request performance. However, the no-LLM
pipeline solves the current 30-case request set faster and without model
calls, so the present experiment does not establish a need for an LLM.
The main lesson is a division of responsibility. The direct-selection
comparison reinforces this conclusion: one LLM sacrifices quality and
variety while the other violates stored dietary conditions despite a
favorable nutrition-fit mean. The four-tool extension adds a second
lesson: all three models detect the injected fault when they reach
validation, but a deterministic release boundary is required to ensure
that the detected fault is actually blocked. Use the model for
language-dependent orchestration only when a harder language benchmark
shows a benefit, and for bounded soft-preference ranking; retain
constraint-sensitive recommendation and identity-binding decisions in
inspectable deterministic components; and evaluate both the trajectory
and the resulting state.

\section{References}\label{references}

{[}1{]} S. Yao, J. Zhao, D. Yu, N. Du, I. Shafran, K. Narasimhan, and Y.
Cao. ``\href{https://openreview.net/forum?id=WE_vluYUL-X}{ReAct:
Synergizing Reasoning and Acting in Language Models}.''
\emph{International Conference on Learning Representations}, 2023.

{[}2{]} Y. Wang, Z. Jiang, Z. Chen, F. Yang, Y. Zhou, E. Cho, X. Fan, Y.
Lu, X. Huang, and Y. Yang.
``\href{https://aclanthology.org/2024.findings-naacl.271/}{RecMind:
Large Language Model Powered Agent for Recommendation}.'' \emph{Findings
of NAACL}, pp. 4351--4364, 2024.
doi:10.18653/v1/2024.findings-naacl.271.

{[}3{]} R. Garimella, B. Srivastava, and A. Sheth.
``\href{https://arxiv.org/abs/2610.23921}{Doing Less with More: A
First-Principles Exploration of the Suitability of Agentic Computing
over Alternative Architectural Choices}.'' 2026.

{[}4{]} S. Yao, N. Shinn, P. Razavi, and K. Narasimhan.
``\href{https://arxiv.org/abs/2406.12045}{\(\tau\)-bench: A Benchmark
for Tool-Agent-User Interaction in Real-World Domains}.'' 2024.

{[}5{]} J. Lu et al.
``\href{https://arxiv.org/abs/2408.04682}{ToolSandbox: A Stateful,
Conversational, Interactive Evaluation Benchmark for LLM Tool Use
Capabilities}.'' 2024.

{[}6{]} S. G. Patil, H. Mao, F. Yan, C. C.-J. Ji, V. Suresh, I. Stoica,
and J. E. Gonzalez.
``\href{https://proceedings.mlr.press/v267/patil25a.html}{The Berkeley
Function Calling Leaderboard (BFCL): From Tool Use to Agentic Evaluation
of Large Language Models}.'' \emph{Proceedings of the 42nd International
Conference on Machine Learning}, pp. 48371--48392, 2025.

{[}7{]} A. Testoni and R. Fernández.
``\href{https://aclanthology.org/2024.eacl-long.16/}{Asking the Right
Question at the Right Time: Human and Model Uncertainty Guidance to Ask
Clarification Questions}.'' \emph{EACL}, pp. 258--275, 2024.
doi:10.18653/v1/2024.eacl-long.16.

{[}8{]} M. J. Q. Zhang and E. Choi.
``\href{https://aclanthology.org/2025.findings-naacl.306/}{Clarify When
Necessary: Resolving Ambiguity Through Interaction with LMs}.''
\emph{Findings of NAACL}, pp. 5541--5558, 2025.
doi:10.18653/v1/2025.findings-naacl.306.

{[}9{]} E. Debenedetti, J. Zhang, M. Balunović, L. Beurer-Kellner, M.
Fischer, and F. Tramèr.
``\href{https://arxiv.org/abs/2406.13352}{AgentDojo: A Dynamic
Environment to Evaluate Prompt Injection Attacks and Defenses for LLM
Agents}.'' 2024.

{[}10{]} E. Debenedetti et al.
``\href{https://arxiv.org/abs/2503.18813}{Defeating Prompt Injections by
Design}.'' 2025.

{[}11{]} V. Nagpal, S. L. Valluru, K. Lakkaraju, and B. Srivastava.
``\href{https://arxiv.org/abs/2406.13714}{BEACON: Balancing Convenience
and Nutrition in Meals With Long-Term Group Recommendations and
Reasoning on Multimodal Recipes}.'' 2024.

{[}12{]} V. Nagpal, S. L. Valluru, K. Lakkaraju, N. Gupta, Z.
Abdulrahman, A. Davison, and B. Srivastava.
``\href{https://arxiv.org/abs/2412.17910}{A Novel Approach to Balance
Convenience and Nutrition in Meals With Long-Term Group Recommendations
and Reasoning on Multimodal Recipes and its Implementation in BEACON}.''
2024.

{[}13{]} Model Context Protocol.
``\href{https://modelcontextprotocol.io/specification/}{Specification}.''
Accessed October 27, 2026.

{[}14{]} LangChain.
``\href{https://docs.langchain.com/oss/python/langchain/agents}{Agents}.''
Accessed October 27, 2026.

{[}15{]} S. Padmanabhuni.
``\href{https://huggingface.co/datasets/srinivaspadmanabhuni/agentic-ai-test-checklist-110}{Agentic
AI Test Checklist 110}.'' Hugging Face dataset, 110 rows, CC BY 4.0.
Accessed November 3, 2026.

\end{document}
